\documentclass[12pt,letterpaper]{article}
\usepackage[T1]{fontenc}
\usepackage{amsmath,amssymb,booktabs,longtable,array,microtype}
\usepackage{newtxtext,newtxmath}
\usepackage[letterpaper,margin=1in]{geometry}
\usepackage{xr}
\externaldocument{zeta_operators_2026}
\usepackage[hidelinks]{hyperref}
\usepackage{xurl}
\title{Revision log: Zeta Functions and Their Associated Operators}
\author{David Joyner\\\normalsize editorial assistance by ChatGPT (OpenAI) and Claude (Anthropic)}
\date{November 28, 2026}
\begin{document}\maketitle
\section{Scope and version}
Current file: \nolinkurl{zeta_operators_2026.tex} (Draft~5).
Draft~4 was produced from \nolinkurl{zeta_operators_corrected_2026c.tex},
renamed \nolinkurl{zeta_operators_corrected_2026d.tex} and then
\nolinkurl{zeta_operators_2026.tex}.  The Draft~4 sections below are
retained as written.
The original 1981 typescript was not inspected in this session.
The revised edition retains existing mathematical arguments,
keeps the theorem labels, and has a newly added elementary
graph-theoretic section. The original claims table is moved
out of the manuscript and reproduced below.

\section{Original-claim status reported by the source draft}
The following is the status assessment from the supplied 2026c
version and should not be mistaken for a new line-by-line
independent verification.


\begin{longtable}{@{}p{0.2\textwidth}p{0.36\textwidth}p{0.36\textwidth}@{}}
\toprule
\textbf{1981 claim} & \textbf{Status} & \textbf{Where in this edition}\\
\midrule
\endhead
Lemma 1(a) & Correct. & Section~\ref{sec:random}.\\
Lemma 1(b) & The constants are not justified. The Bernstein inequality is
misnormalized. & Replaced by an almost-sure bound from Hoeffding's
inequality.\\
Lemma 1(c) & The conclusion is true. The counting argument does not
prove it. & New proof by conditioning, Theorem~\ref{thm:random}.\\
Lemma 1(d) & The inversion step is insufficient. & Replaced by an
almost-sure inversion.\\
Theorem 1 & True. The proof has a wrong von Mangoldt weight. &
Theorem~\ref{thm:random}.\\
Theorem 2 & True. The last step of the proof is false. &
Theorem~\ref{thm:sparse}. A $\theta$-free form and an equivalence with RH
are new (Corollaries~\ref{cor:thetafree} and~\ref{cor:RHequiv}).\\
Theorem 3 & The generator $a_1=1$ is impossible. Otherwise true after
repair. & Corollary~\ref{cor:sparseA}, from the comparison principle
(Theorem~\ref{thm:comparison}).\\
Corollary 1 & True. & Corollary~\ref{cor:pointwise}.\\
Lemma 2 & True. The proof is a sketch. & Lemma~\ref{lem:density}.\\
Lemma 3 & True. The density proof needs repair. &
Lemma~\ref{lem:density}.\\
Theorem 4(b) & True. & Theorem~\ref{thm:spectral}.\\
Theorem 4(a) & False. & Theorem~\ref{thm:manyeigen} gives
counterexamples. The correct statement is a joint-kernel statement,
Theorem~\ref{thm:spectral}.\\
Remark after Theorem 4 & True and incomplete. Complex eigenvalues
exist. & Theorem~\ref{thm:manyeigen}. Every $\lambda\notin[0,\infty)$ is an
eigenvalue.\\
Theorem 5 & Unproved. Where it makes sense it is vacuous. Both operators
then have every negative number as an eigenvalue. &
Remark~\ref{rem:thm5}. Replaced by Proposition~\ref{prop:EA}.\\
Theorem 6 & False where the operator is defined. &
Corollary~\ref{cor:thm6false}. Replaced by Theorem~\ref{thm:slit} and
Proposition~\ref{prop:KU}.\\
Closing remarks on pair correlation & Speculative. &
Section~\ref{sec:conclusion}.\\
\bottomrule
\end{longtable}


\section{Graph additions}
The new results consist of a discrete-Stokes identity on the
prime path; a weighted-boundary estimate for comparison
operators on $S_k$; the generalized-integer exponent-vector
graph; and periodic Fourier weights on directed cycles.
The graph-boundary theorem is proved under an explicit
summability hypothesis. No graph-Laplacian isospectrality
or full-operator similarity is claimed.

\section{Style and verification}
Retained the original \texttt{article} class, changed
to 12-point Times, US Letter, one-inch margins, and
inserted an introduction and pre-1981 literature review.
Placed the literature review and notation within Section~1, and retained the status table here.
Citations were re-keyed to author initials and years.
The original analytic proofs were not independently
re-refereed in full. The Grosswald--Schnitzer citation
was checked in the 1978 \emph{Pacific Journal of Mathematics}
original. Other bibliographic details were inherited
from the attached source and require independent checking.
AI review is not independent certification.

\section{November 28, 2026: RH editorial corrections}
The previous revision omitted the November 2026 RH formalization from
the historical paper's post-1981 literature review, although a separate
modern companion note discussed it. This pass adds citations to the
joint statement of the Lean FRO and Mathlib Initiative and to
Rubinstein's independent checks. It annotates von Koch's historically
conditional estimate, the lemma on $p_n-u_n$, the sparse-exception
equivalence corollary, and the $S_1$ algebraic joint-kernel corollary.

The mathematical statements and proofs retain their historical
logical form. The new editorial comments concern the status of RH
as of November 28, 2026; they are not an independent check of the
formal proof by ChatGPT, nor an independent re-refereeing of the
existing operator-theoretic arguments.

Sources consulted for status:
\url{https://leanprover-community.github.io/blog/posts/riemann-hypothesis/}
and
\url{https://math.uwaterloo.ca/~mrubinst/alpoge_claude_grh/}.

\section{Draft 4 summary (November 28, 2026)}
Moved the literature review and notation under the Introduction; added
one sentence to each existing glossary entry and 22 further technical
terms; preserved the historical 1981 mathematical claims.  The RH status
and conditional operator implications are in the separate modern
companion \nolinkurl{rh_2026_operator_implications.tex}.  (This section
replaces the former file \nolinkurl{revision_log_draft4.md}; the former
file \nolinkurl{rh_editorial_revision_2026_11_28.md} is summarized in the
preceding section.)

\section{Draft 5 (November 28, 2026), prepared by Claude}
Claude (Anthropic) refereed Draft~4 together with the companion papers
in \texttt{grosswald-schnitzer-graph} and the RH companion note, and
made the following changes.
\begin{itemize}
\item Added Remark~\ref{rem:synthesis-2026}.  The companion paper
  \emph{Spectral invariance under Grosswald--Schnitzer prime
  perturbations} proves the closed joint-kernel statement asked for in
  Open Problem~\ref{prob:synthesis}.  The remark records this and the
  resulting closed form of the RH equivalence for $k=1$.  The proof in
  the companion paper was checked by Claude.  Its Remark~2.4 records
  that it applies to the spaces of this edition, which use $N$
  derivatives instead of $k$.
\item The editorial note in the Introduction and the corollary after
  Theorem~\ref{thm:betafree} now point to that remark.  They no longer
  say that spectral synthesis is open.
\item The RH result is attributed as \emph{found with Anthropic's
  Claude}, citing Rubinstein's page, since the Lean FRO and Mathlib
  statement does not mention Claude.  The title of the Rubinstein
  citation was corrected to the page's actual title, and an access date
  was added.
\item Glossary: merged the near-duplicates.  ``Discrete divergence'' is
  now folded into ``Directed incidence boundary''.  ``Graph boundary
  variation'' and ``Weighted total variation'' are replaced by
  ``Weighted boundary variation'', whose former pointer to a theorem is
  replaced by a self-contained definition.  The glossary was
  alphabetized, and the entries ``Jordan block'' and ``Resolvent set''
  were corrected.
\item Introduced section-based equation numbering, following the
  style guide.  Theorem numbering is unchanged.
\item Rewrote the Acknowledgments as a list of the five 2026 editions
  and their preparers.
\item Fixed this log's \verb|\externaldocument| target, which still named
  the pre-rename file, so all cross-references resolve again.
\end{itemize}
Not changed: the discussion in the body of the errors in the 1981
typescript, including the counterexamples to the original Theorems~4(a)
and~6, is kept.  This is a deliberate exception to the style guide's
preference for moving error discussion into the revision log: in a
corrected historical edition, the counterexamples are part of the
content.  The 1981-era proofs of Sections~2--9 were checked by Claude in
the referee pass; no mathematical errors were found.

\section{AI contributions}
The source credits ChatGPT (OpenAI) and Claude (Anthropic)
for the earlier corrected editions. ChatGPT (OpenAI)
supplied graph proofs, examples, structure and editing
in this new draft. Modern results are not used as inputs to the historical proofs; 2026 editorial citations distinguish formal RH verification from the original conditional statements.
\end{document}
